\subsection[Binary Insertion Sort]{Binary Insertion Sort}
    \subsubsection{Overview}
        \begin{itemize}
            \item \textbf{History and Origin:} Binary insertion sort is a well-known variant of insertion sort that uses binary search to determine the correct insertion position. This technique was mentioned by John Mauchly in 1946 in one of the earliest published discussions on computer sorting.
            \item \textbf{Efficiency:} Best for small datasets or nearly sorted/sorted data. Worst for reverse-sorted datasets.
            \item \textbf{Application:} Instead of using insertion sort when the subarray size becomes small in hybrid algorithms such as Timsort, Introsort, or Quicksort, binary insertion sort can be used for better efficiency. It is also useful in problems where the cost of comparisons is relatively high.
        \end{itemize}
    \subsubsection{Idea}
        \begin{itemize}
            \item Let the list be divided by sub-lists, separating them into sorted and unsorted sections.
            \item Take the first element of the unsorted list and insert it into its proper place in the sorted one.
            \item After each placement, the size of the sorted region grows by 1 and the size of the unsorted region shrinks by 1.
        \end{itemize}

    \subsubsection{Step-by-Step Algorithm}
        \begin{itemize}
            \item Binary insertion sort uses the same code structure as Insertion Sort, except for the process of finding the correct position. Therefore, only this step is described below.
            \item Find the correct position \textit{pos} in $a[0\ldots (i - 1)]$ to insert a[i], where a[pos - 1] $\leq$ a[i] < a[pos].
            \item \textbf{Step 1:} Set val = a[i], \textit{L} = 0, \textit{R} = i - 1 and \textit{pos} = i.
            \item \textbf{Step 2:} While L $\leq$ R.
                \begin{itemize}
                    \item Set Mid = (L + R) / 2.
                    \item If a[Mid] $\geq$ val then set R = Mid - 1 and pos = Mid.
                    \item Else set L = Mid + 1.
                \end{itemize}
            \item \textbf{Step 3:} Move forward a[pos \ldots (i - 1)] to one element. After that, set a[pos] = val. 
        \end{itemize}

    \subsubsection{Pseudo Code}
        \begin{itemize}
            \item \textbf{Algorithm:} Binary Insertion Sort
            \item \textbf{Input:} An array A of n elements
        \end{itemize}
        \begin{verbatim}
        1. for i ← 1 to length(A)-1
        2.      value ← A[i]
        3.      l ← 0, r ← i-1, pos ← i
        4.      while l <= r
        5.           m ← l + (r-l) / 2
        6.           if A[m] >= value
        7.                pos ← m
        8.                r ← m-1
        9.           else
        10.               l ← m+1
        11.     for j ← i downto pos+1
        12.          A[j] ← A[j-1]
        13.     A[pos] ← value
        \end{verbatim}
    \subsubsection{Complexity}
        \begin{itemize}
            \item \textbf{Time Complexity:}
            \begin{itemize}
                \item Finding correct position takes $O(logn)$ but the process of insertion still takes $O(n)$.
                \item Best case: $O(nlogn)$
                    \subitem When the input is already sorted.
                \item Average case: $O(n^2)$
                    \subitem When the input is random.
                \item Worst case: $O(n^2)$ 
                    \subitem When the input is reverse-sorted.
            \end{itemize}
            \item \textbf{Space Complexity:} $O(1)$ 
                \subitem In-place algorithm. Requires only a constant amount of extra memory.
        \end{itemize}
        
    \subsubsection{Variants and Optimizations}
        \begin{itemize}
            \item \textbf{Binary Insertion Sort:}
            \begin{itemize}
                \item Use Binary Search to find the position of an element within the sorted sub-array which is greater than or equal to A[i].
                \item While Binary Search reduces the number of comparisons to $O(n \log n)$, the complexity of data movement remains $O(n^2)$ because elements must still be shifted to make space for the new entry.
            \end{itemize}
        \end{itemize}